% Part: proof-theory
% Chapter: sequent-calculus
% Section: introduction

\documentclass[../../../include/open-logic-section]{subfiles}

\begin{document}

\olfileid{pt}{seq}{int}

\olsection{పరిచయం}

సీక్వెంట్ కలనశాస్త్రాలు 1930లలో గెర్హార్డ్ గెంట్సెన్ ప్రవేశపెట్టిన \tetoken{నిరూపణ}{!!{proof}} వ్యవస్థలు. వాటి ప్రధాన ఉద్దేశం విషయాలను \tetoken{నిరూపించడానికి}{!!{proving}} ఆచరణాత్మక పద్ధతి కావడం కాదు; \tetoken{నిరూపణల}{!!{proof}s} సాధ్యమైన నిర్మాణం, వాటి లక్షణాలపై కొన్ని ఫలితాలను నిరూపించడానికి గెంట్సెన్ వాటిని వాడారు. ఇలాంటి ప్రశ్నలే నిరూపణ సిద్ధాంతం పరిశీలిస్తుంది.

పేరు సూచించినట్లే, సీక్వెంట్ కలనశాస్త్రాలు \tetoken{సూత్రాలపై}{!!{formula}s} కాక \emph{సీక్వెంట్లపై} పనిచేస్తాయి. సీక్వెంట్ అనేది \tetoken{సూత్రాల}{!!{formula}s} వరుసలు లేదా బహుసమితుల జత. గెంట్సెన్ అసలు వ్యవస్థలు (\Log{LK}, \Log{LJ}) వరుసలను వాడాయి; ఇప్పుడు నిరూపణ సిద్ధాంతవేత్తలు బహుసమితులను ఇష్టపడతారు. బహుసమితి సమితి వంటిదే; కానీ ఒక \tetoken{మూలకాన్ని}{!!{element}} ఒకటి కంటే ఎక్కువసార్లు కలిగి ఉండవచ్చు. సమితుల్లోలాగే \tetoken{మూలకాల}{!!{element}s} క్రమానికి ప్రాధాన్యం లేదు. కాబట్టి $\{!A, !A, !B\}$, $\{!A, !B, !A\}$ ఒకే బహుసమితిని వర్ణిస్తాయి. కానీ అది $\{!A,
!B\}$, $\{!A, !A, !B, !B\}$ల నుంచి భిన్నం; వాటిలో $!A$, $!B$ల ప్రతుల సంఖ్యలు వేరు.

నిరూపణ సిద్ధాంతంలో \tetoken{సూత్రాల}{!!{formula}s} బహుసమితులకు సంప్రదాయంగా పెద్ద గ్రీకు అక్షరాలను వాడతారు. కాబట్టి $\Gamma$, $\Delta$, $\Pi$, $\Lambda$, $\Theta$ అని రాసినప్పుడు, పరిశీలిస్తున్న తర్కంలోని \tetoken{సూత్రాల}{!!{formula}s} బహుసమితులని ఉద్దేశిస్తాం. అలాగే $\tuple{\Gamma, \Delta}$ అని రాయడానికి బదులు, రెండు భాగాలను సీక్వెంట్ గుర్తుతో వేరు చేస్తారు; మనం~$\Sequent$ వాడతాం.

\begin{defn}[సీక్వెంట్]
\emph{సీక్వెంట్} అనేది
\[
\Gamma \Sequent \Delta
\]
రూపం గల వ్యక్తీకరణ; ఇక్కడ $\Gamma$, $\Delta$ పరిమిత (ఖాళీ కావచ్చు) \tetoken{సూత్రాల}{!!{formula}s} బహుసమితులు. $\Gamma$ను \emph{పూర్వపక్షం}, $\Delta$ను \emph{ఉత్తరపక్షం} అంటారు.
\end{defn}

$\Gamma$లోని \tetoken{సూత్రాల}{!!{formula}s} సంయోగం~$!G$ అనుకుందాం (ఖాళీ సంయోగాన్ని~$\ltrue$గా తీసుకుంటాం); $\Delta$లోని \tetoken{సూత్రాల}{!!{formula}s} వియోగం~$!D$ అనుకుందాం (ఖాళీ వియోగాన్ని~$\lfalse$గా తీసుకుంటాం). అప్పుడు ఏదో నిర్మాణం~$\Struct{M}$లో $!G \lif !D$ సత్యమైతేనే $\Gamma \Sequent \Delta$ సత్యం. సాంప్రదాయ తర్కంలో దీని అర్థం: $\Gamma$లో కనీసం ఒక \tetoken{సూత్రం}{!!{formula}s} అసత్యం లేదా $\Delta$లో కనీసం ఒక \tetoken{సూత్రం}{!!{formula}s} సత్యం.

$\Gamma$ \tetoken{సూత్రాల}{!!{formula}s} బహుసమితి అయితే, $\Gamma$కు $!A$ను చేర్చిన ఫలితాన్ని $\Gamma, !A$ (లేదా $!A, \Gamma$) అని రాస్తాం. $\Delta$ కూడా \tetoken{సూత్రాల}{!!{formula}s} బహుసమితి అయితే, $\Gamma, \Delta$ వాటి బహుసమితి సమ్మేళనం: $\Gamma$లో $!A$ \emph{బాహుళ్యం}~nతో ఉంటే (అంటే~$!A$కు $n$ ప్రతులు ఉంటే), $\Delta$లో బాహుళ్యం~$m$తో ఉంటే, $\Gamma, \Delta$లో $!A$ బాహుళ్యం~$n+m$. సీక్వెంట్‌లో ఒకవైపు బహుసమితి ఖాళీ అయితే సాధారణంగా ఆ వైపును ఖాళీగా వదిలేస్తాం. ఉదాహరణకు $\Sequent !A$ పూర్వపక్షం ఖాళీ, ఉత్తరపక్షంలో $!A$ బాహుళ్యం~$1$ ఉన్న సీక్వెంట్.

సీక్వెంట్ కలనశాస్త్రంలో \tetoken{ఒక నిరూపణ}{!!^a{proof}} సీక్వెంట్ల వృక్షం. పై (ప్రారంభ) సీక్వెంట్లు పరిశీలిస్తున్న వ్యవస్థలో అక్షయాలు కావాలి; ఒక సీక్వెంట్ మరో ఒకటి లేదా రెండు సీక్వెంట్ల కింద ఉంటే, అది ఆ వ్యవస్థ నిగమన నియమంతో సరిగ్గా అనుసరించాలి. ముందుగా సాంప్రదాయ తర్కానికి రెండు వేర్వేరు వ్యవస్థలు $\Log{G1c}$, $\Log{G3c}$లను పరిశీలిస్తాం. వాటి అక్షయాలు, నియమాలు \cref{pt:seq:int:tab:G1c,pt:seq:int:tab:G3c}లో ఉన్నాయి. \oliflabeldef{fol:seq::chap}{ (గెంట్సెన్ అసలు వ్యవస్థ \Log{LK}ను \olref[fol][seq][]{chap}లో వర్ణించాం.)}{}

\subfile{rules-G1c}
\subfile{rules-G3c}

ఈ వ్యవస్థల్లో సూక్ష్మ తేడాలు ఉన్నాయి; వాటివల్ల నిరూపణ-సిద్ధాంత లక్షణాలూ వేరుగా ఉంటాయి. \Log{G1c} అక్షయాలు $!A
\Sequent !A$ రూపంలోనే ఉంటాయి; ఇతర \tetoken{సూత్రాలు}{!!{formula}s} అనుమతించబడవు. \Log{G3c} అక్షయాలు $!C, \Gamma \Sequent \Delta, !C$ రూపంలో ఉంటాయి; $\Gamma$, $\Delta$లలో ఏవైనా “పక్క” \tetoken{సూత్రాలు}{!!{formula}s} ఉండవచ్చు, కానీ రెండు వైపులా వచ్చే \tetoken{సూత్రం}{!!{formula}}~$!C$ \emph{పరమాణువు} కావాలి. \Log{G1c}లో $\LeftR{\land}$, $\RightR{\lor}$లకు రెండేసి రూపాలు; \Log{G3c}లో ఒక్కొక్కటే. ఇంకా \Log{G1c}లో “సంకోచన” (\LeftR{\Contraction}, \RightR{\Contraction}), “బలహీనీకరణ” (\LeftR{\Weakening}, \RightR{\Weakening}) అనే అదనపు “నిర్మాణ నియమాలు” ఉన్నాయి.

\Log{G1c}, \Log{G3c} సాంప్రదాయ తర్కానికి సమర్థమైనవి, పరిపూర్ణమైనవి. అంటే వీటిలో \tetoken{నిరూప్యమైన}{!!{provable}} ప్రతి సీక్వెంట్ ప్రతి \tetoken{నిర్మాణంలో}{!!{structure}} సత్యం; ప్రతి \tetoken{నిర్మాణంలో}{!!{structure}} సత్యమైన ప్రతి సీక్వెంట్‌కు \tetoken{ఒక నిరూపణ}{!!a{proof}} ఉంది. కాబట్టి సీక్వెంట్‌కు \tetoken{ఒక నిరూపణ}{!!a{proof}} \emph{ఉందా} అనే ప్రశ్నకు ఇతర తార్కిక పద్ధతులతోనూ (ఉదా., నమూనా సిద్ధాంతంతో) సమాధానం ఇవ్వవచ్చు. రెండు వ్యవస్థలూ ప్రతి \tetoken{నిర్మాణంలో}{!!{structure}} సత్యమైన సీక్వెంట్లనే \tetoken{నిరూపిస్తాయి}{!!{prove}} కనుక, అవి ఒకే సీక్వెంట్లను నిరూపిస్తాయి.

కానీ నిరూపణ సిద్ధాంతం ఏదైనా నిరూపించవచ్చా అన్నదే కాక, \emph{ఎలా} నిరూపించవచ్చో కూడా పరిశీలిస్తుంది. “ఒక నియమాన్ని ఎప్పుడూ నివారించవచ్చా?”, “కొత్త సీక్వెంట్లు నిరూప్యమవకుండా ఈ నియమాన్ని వ్యవస్థకు చేర్చవచ్చా?”, “ఒక వ్యవస్థలోని \tetoken{నిరూపణలను}{!!{proof}s} మరో వ్యవస్థలోని \tetoken{నిరూపణలుగా}{!!{proof}s} మార్చవచ్చా?” వంటి ప్రశ్నలను సిద్ధాంతవేత్తలు అడుగుతారు. సమాధానం “అవును” అయితే, దాని నిరూపణ తరచుగా \tetoken{నిరూపణల}{!!{proof}s} సంక్లిష్టతపై ఆగమనంతో లేదా సమానంగా \tetoken{నిరూపణల}{!!{proof}s} నిర్మాణంపై ఆగమనంతో ఇస్తారు. ఇది వాస్తవానికి నిరూపణనే కాక (ఉదా., \Log{G1c} ఒక సీక్వెంట్‌ను నిరూపిస్తే \Log{G3c} కూడా దాన్ని నిరూపిస్తుంది), ఒక విధానాన్నీ ఇస్తుంది (ఉదా., \Log{G1c}లోని \tetoken{నిరూపణలను}{!!{proof}s} \Log{G3c}లోని \tetoken{నిరూపణలుగా}{!!{proof}s} మార్చే విధానం).

నిరూపణ సిద్ధాంతంలో కేంద్రీయ ఫలితం \emph{కట్-తొలగింపు సిద్ధాంతం}. కట్ నియమం
\[
\Axiom$\Gamma \fCenter \Delta, !A$
\Axiom$!A, \Pi \fCenter \Lambda$
\RightLabel{\Cut}
\BinaryInf$\Gamma, \Pi \fCenter \Delta, \Lambda$
\DisplayProof
\]
సీక్వెంట్ వ్యవస్థలకు చేర్చినా, ఏ సీక్వెంట్లు \tetoken{నిరూప్యమైనవో}{!!{provable}} మారదని అది చూపుతుంది. అంతేకాక \Log{G1c} (లేదా \Log{G3c}) నియమాలతోపాటు \Cut{}ను వాడే ఏ \tetoken{నిరూపణనైనా}{!!{proof}}, \Log{G1c} (లేదా \Log{G3c}) నియమాలనే వాడే నిరూపణగా మార్చే విధానాన్ని ఇస్తుంది. అంటే \tetoken{నిరూపణలలోని}{!!{proof}s} \Cut{} నియమాన్ని ఈ విధానం “తొలగిస్తుంది”; అందుకే ఆ పేరు.

\end{document}
